\documentclass[11pt]{article}
\usepackage[letterpaper,margin=1in]{geometry}
\usepackage{amsmath,amssymb,amsthm,mathtools}
\usepackage{booktabs,tabularx,array,microtype,xurl}
\usepackage[colorlinks=true,linkcolor=blue,citecolor=blue,urlcolor=blue]{hyperref}
\newtheorem{theorem}{Theorem}[section]
\newtheorem{proposition}[theorem]{Proposition}
\newtheorem{lemma}[theorem]{Lemma}
\newtheorem{corollary}[theorem]{Corollary}
\theoremstyle{remark}
\newtheorem{remark}[theorem]{Remark}
\newcommand{\R}{\mathbb R}
\newcommand{\C}{\mathbb C}
\newcommand{\E}{\mathbb E}
\newcommand{\dd}{\,\mathrm d}
\newcommand{\Var}{\operatorname{Var}}
\newcommand{\Cov}{\operatorname{Cov}}
\newcommand{\LP}{\mathcal{LP}}
\newcommand{\norm}[1]{\left\lVert#1\right\rVert}
\numberwithin{equation}{section}
\title{Compatibility obstructions for constructive stable theta models\\
\large Independent seeds and ferromagnetic quartic networks}
\author{Drafted for Edward Baker}
\date{17 October 2026\\\small Unreviewed working draft}
\begin{document}
\maketitle
\begin{abstract}
We examine explicit real-zero constructions as possible approximants to the
complete theta heat flow. A normalized ferromagnetic comparison shows that
any quartic network approximating theta must lose every persistent diffuse
quartic site at its readout scale, even when the number of sites and all
other couplings vary without bounds. Independently scaled pure-quartic
convolutions, with arbitrary finite counts and an additional Gaussian
component, cannot approximate theta even weakly: field growth forces
infinitesimal weights, after which the limit must be Gaussian. We also collect
finite interval obstructions for independent Bernoulli products, single
quartic laws, and a fixed Bessel base with nonnegative even multipliers.
Classical Lee--Yang theory supplies explicit coupled-quartic seeds and
their Ising realizations; a certified negative sixth potential coefficient
shows why a local source exclusion does not extend to every such marginal.
No theta approximation theorem or conclusion about the Riemann hypothesis
is obtained. Infinitesimal interacting arrays and degenerating elementary
measures remain outside the exclusions.
\end{abstract}

\noindent\textbf{Draft status.}
This manuscript was requested before formal round review. It integrates
round-one Notes 01--03 and their small replay records. Derivations,
certificates and authoring checks used \texttt{gpt-6.1-sol} (Codex), reasoning
effort \texttt{ultra}. The checks are internal and same-model; cross-model
round review and specialist human review remain pending. No literature
priority is claimed.

\section{Target, conventions and scope}\label{sec:setup}

Use the complete theta source and heat transform
\begin{align}
\Phi(u)&=\sum_{n\ge1}
 (2\pi^2n^4e^{9u}-3\pi n^2e^{5u})e^{-\pi n^2e^{4u}},
 \qquad u\ge0,\label{eq:phi}\\
H_t(z)&=\int_0^\infty e^{tu^2}\Phi(u)\cos(zu)\dd u,
 \qquad H_0(z)=\tfrac18\xi((1+iz)/2).\label{eq:heat}
\end{align}
The source extends evenly to the real line. The xi identity and the heat
normalization are classical; see Polymath \cite{polymath}. In particular
\(\partial_tH_t=-\partial_z^2H_t\). Set
\[
 I_{2k}=\int_0^\infty u^{2k}\Phi(u)\dd u,\qquad
 m_{2k}=I_{2k}/I_0,\qquad
 M_\theta(h)=H_0(-ih)/I_0=\frac{\xi((1+h)/2)}{\xi(1/2)}.
\]
Thus \(M_\theta\) is the normalized Laplace transform of the probability
density \(\Phi(u)/(2I_0)\) on \(\R\). All finite complex-field and heat
derivatives converge absolutely on compact sets.

The inherited outward moment enclosure \cite{scout} gives
\begin{equation}\label{eq:moments}
 0.01155234249<m_2<0.01155265063,\qquad
 2.79098616<\frac{m_4}{m_2^2}<2.79121956.
\end{equation}
The strict consequence
\(\kappa_4(M_\theta)=m_4-3m_2^2<0\) supplies an optional check below;
the new exclusion theorems have analytic proofs without this enclosure.
These are source-moment enclosures, not zero certificates.

The other growth input is elementary and unconditional.
\begin{lemma}[Theta field growth]\label{lem:theta}
As \(h\to+\infty\),
\begin{equation}\label{eq:theta-growth}
\log M_\theta(h)
 =\frac h4\log h-\frac h4(1+\log(4\pi))
  +\frac74\log h+O(1).
\end{equation}
In particular \(\log M_\theta(h)=o(h^{4/3})\).
\end{lemma}
\begin{proof}
Insert \(s=(1+h)/2\) into
\(\xi(s)=s(s-1)\pi^{-s/2}\Gamma(s/2)\zeta(s)/2\).
Stirling's formula \cite[5.11.1]{dlmf}, with \(s/2=(h+1)/4\), and the
Dirichlet-series estimate \(\zeta(s)=1+O(2^{-s})\) give the expansion.
The \(\log h\) coefficient is \(2-1/4=7/4\).
\end{proof}

We use the term stable for the real-zero Fourier or imaginary-zero Laplace
property of the specified entire transform. Simplicity is not included.
A construction of stable approximants converging to \(H_0\) on every
complex compact would imply RH by Hurwitz's theorem. Our results are
scoped obstructions to particular constructions, rather than evidence
from finite checks for that global objective.

\section{Ferromagnetic comparison and persistent quartic sites}
\label{sec:persistent}

Consider a finite probability law
\begin{equation}\label{eq:network}
\dd\nu(\mathbf x)=Z^{-1}\exp\left\{
 -\sum_{j=1}^d\alpha_jx_j^4-\sum_j\beta_jx_j^2+
 \sum_{i<j}K_{ij}x_ix_j\right\}\dd\mathbf x,
\quad \alpha_j>0,\quad K_{ij}\ge0,
\end{equation}
with arbitrary real \(\beta_j\), and readout \(S=\sum_jb_jx_j\), \(b_j>0\).
Its normalized transform is \(M(h)=\E_\nu e^{hS}\). Quartic tails justify
every finite derivative used in this section.

Absorb the readout into the coordinates:
\begin{equation}\label{eq:effective}
 Y_j=b_jx_j,\qquad
 a_j=\alpha_j/b_j^4,\quad c_j=\beta_j/b_j^2,\quad
 L_{ij}=K_{ij}/(b_ib_j).
\end{equation}
Write
\begin{equation}\label{eq:site}
 m_{a,c}(h)=
 \frac{\int_\R e^{-ay^4-cy^2+hy}\dd y}
      {\int_\R e^{-ay^4-cy^2}\dd y},
 \qquad a>0,\quad c\in\R.
\end{equation}

\begin{lemma}[Normalized ferromagnetic comparison]\label{lem:gks}
For every real \(h\ge0\), a network \eqref{eq:network} satisfies
\begin{equation}\label{eq:comparison}
 M(h)\ge\prod_{j=1}^d m_{a_j,c_j}(h)\ge m_{a_k,c_k}(h)
 \qquad(1\le k\le d).
\end{equation}
There is no bound on \(d\), the other parameters, or the coupling sizes.
\end{lemma}
\begin{proof}
We give the classical GKS duplication argument in the precise continuous
class used here. For any two spin monomials \(Y_A,Y_B\), including repeated
indices, the covariance is nonnegative under nonnegative fields and pair
couplings. Indeed, take two copies \(y,z\) and put
\(u=(y+z)/2\), \(v=(y-z)/2\).
Twice the covariance is the expectation of
\((y_A-z_A)(y_B-z_B)\) under the product law.
Each difference has nonnegative polynomial coefficients in \(u,v\):
exactly the terms with an odd total number of \(v\)'s survive.
The doubled interaction and field are
\[
 2\sum_{i<j}L_{ij}(u_iu_j+v_iv_j)+2\sum_jh_ju_j.
\]
Their exponential has nonnegative coefficients. Keep the doubled
single-site density unexpanded. It is invariant independently under
each \(u_j\mapsto-u_j\) and \(v_j\mapsto-v_j\), by evenness and exchange
of the two copies. Each integrated monomial is therefore zero or
nonnegative. Quartic tails justify absolute integration of the expanded
interaction and polynomial factors. This proves the covariance inequality.

Let \(\langle\cdot\rangle_h\) denote the common-field Gibbs law. Then
\[
 \frac{\dd}{\dd h}\langle Y_iY_j\rangle_h
  =\sum_k\Cov_h(Y_iY_j,Y_k)\ge0.
\]
For the normalized transform,
\[
 \partial_{L_{ij}}\log M(h)
  =\langle Y_iY_j\rangle_h-\langle Y_iY_j\rangle_0\ge0.
\]
Increase the couplings from zero to \(L\). This yields the product
comparison. Each even single-site transform is an expectation of a
cosh and is at least one. The proof concerns normalized transforms;
dropping couplings from an unnormalized partition function would not
establish the assertion.
\end{proof}

\begin{lemma}[Single-site monotonicity and growth]\label{lem:site}
For \(h\ge0\), \(m_{a,c}(h)\) is nonincreasing in \(a\) and \(c\).
For fixed \(A>0,B\ge0\),
\begin{equation}\label{eq:quartic-growth}
 \log m_{A,B}(h)
 =\frac{3}{4^{4/3}}A^{-1/3}h^{4/3}
   +O(h^{2/3}+\log h).
\end{equation}
An explicit lower bound, valid for \(h\ge1\), is
\begin{align}
 \log m_{A,B}(h)&\ge F_{A,B}(h),\label{eq:field-lower}\\
 F_{A,B}(h)&=hu-A(u+h^{-1})^4-B(u+h^{-1})^2-\log h
  -\log\frac{\Gamma(1/4)}{2A^{1/4}},
 \quad u=(h/(4A))^{1/3}.\notag
\end{align}
\end{lemma}
\begin{proof}
Differentiating the normalized expectation of \(\cosh(hY)\) gives
\[
 \partial_a m_{a,c}(h)=-\Cov(\cosh(hY),Y^4),\qquad
 \partial_c m_{a,c}(h)=-\Cov(\cosh(hY),Y^2).
\]
Both covariances are nonnegative: each pair consists of increasing
functions of \(Y^2\), as is seen from the independent-copy covariance
identity. For \eqref{eq:field-lower}, integrate the numerator over
\([u,u+h^{-1}]\), and bound the denominator by
\(\int_\R e^{-Ay^4}\dd y=\Gamma(1/4)/(2A^{1/4})\).
The leading term is the positive term in \eqref{eq:quartic-growth}.
The matching upper asymptotic follows by scaling \(y=h^{1/3}v\) and
maximizing \(v-Av^4\); the quadratic term contributes \(O(h^{2/3})\).
\end{proof}

\begin{theorem}[No persistent diffuse quartic site]\label{thm:persistent}
Let \(M_n\) be normalized transforms of any sequence of networks
\eqref{eq:network}, with dimensions and couplings allowed to vary without
bounds. If \(M_n(h)\to M_\theta(h)\) for every real \(h\ge0\), then for
every \(A>0,B\ge0\), eventually no site satisfies both
\begin{equation}\label{eq:rectangle}
 \alpha_{n,j}/b_{n,j}^4\le A,\qquad
 \beta_{n,j}/b_{n,j}^2\le B.
\end{equation}
In particular a fixed nontrivial quartic site cannot be retained while
other ferromagnetic sites are appended to approximate theta.
\end{theorem}
\begin{proof}
If such a site persists along a subsequence, Lemmas \ref{lem:gks} and
\ref{lem:site} give \(M_n(h)\ge m_{A,B}(h)\) along it, for every \(h\ge0\).
Equations \eqref{eq:theta-growth} and \eqref{eq:field-lower} supply one
fixed finite \(h_*\) with \(m_{A,B}(h_*)>M_\theta(h_*)\), contradicting
convergence at that field.
\end{proof}

\begin{corollary}[Infinitesimal readouts on bounded site shapes]
\label{cor:infinitesimal}
Put
\(\ell_{n,j}=b_{n,j}\alpha_{n,j}^{-1/4}\) and
\(r_{n,j}=\beta_{n,j}\alpha_{n,j}^{-1/2}\).
Under the hypotheses of Theorem \ref{thm:persistent}, for every fixed
real \(R\),
\begin{equation}\label{eq:small-site}
 \max_{j:r_{n,j}\le R}\ell_{n,j}\longrightarrow0,
\end{equation}
with the empty maximum defined as zero.
\end{corollary}
\begin{proof}
If \(r\le R\) and \(\ell\ge\varepsilon\), then
\(a=\ell^{-4}\le\varepsilon^{-4}\) and
\(c=r/\ell^2\le\max(R,0)/\varepsilon^2\), contradicting
\eqref{eq:rectangle}.
\end{proof}

\begin{remark}[Finite-field witness and limitations]\label{rem:gap}
Whenever a chosen \(h_*\ge1\) satisfies
\(F_{A,B}(h_*)>\log M_\theta(h_*)\), every network retaining a site
in \eqref{eq:rectangle} has
\[
 |M(h_*)-M_\theta(h_*)|
 \ge e^{F_{A,B}(h_*)}-M_\theta(h_*)>0.
\]
No numerical \(h_*\) is claimed. The theorem needs convergence of
exponential observables. Without an additional uniform-integrability
argument, weak convergence does not permit this theorem to be applied.
It does not exclude infinitesimal interacting arrays or discrete-site
degenerations. Quartic field growth must not be ascribed to every Ising
scaling limit; for example, critical planar limits have different tail
exponents \cite{cgn}.
\end{remark}

\section{Independent pure-quartic convolutions}\label{sec:indquartic}

Let \(X\) have normalized density proportional to \(e^{-x^4}\), and set
\(m(h)=\E e^{hX}\). Consider all finite products
\begin{equation}\label{eq:pure-products}
 M_n(h)=e^{\sigma_n^2h^2/2}\prod_{j=1}^{d_n}m(b_{n,j}h),
 \qquad \sigma_n^2\ge0,\quad b_{n,j}\ge0.
\end{equation}
The counts and coefficients are unrestricted; the maximum scale in an
empty row is defined to be zero. These represent independent
scaled pure-quartic blocks plus an independent centered Gaussian.
Each construction is classically stable \cite{sg,fr}.

\begin{theorem}\label{thm:pure}
The transforms \eqref{eq:pure-products} cannot converge to \(M_\theta\)
uniformly on every complex compact. This conclusion has an entirely
analytic proof and does not require a numerical theta moment enclosure.
\end{theorem}
\begin{proof}
If \(\max_jb_{n,j}\ge\varepsilon>0\) along a subsequence, evenness gives
\(M_n(h)\ge m(\varepsilon h)\) for every real \(h\ge0\). Quartic growth
\eqref{eq:quartic-growth} contradicts \eqref{eq:theta-growth} at one
sufficiently large fixed field. Hence \(\max_jb_{n,j}\to0\).

Derivative convergence at zero bounds the variances. With \(v=\E X^2>0\),
\[
 \sigma_n^2+v\sum_jb_{n,j}^2=\Var(S_n)\le C.
\]
For fixed real \(h\), the analytic expansion at zero is
\(\log m(bh)=vb^2h^2/2+O_h(b^4)\), uniformly once the maximum scale is
small. Since
\(\sum_jb_{n,j}^4\le(\max_jb_{n,j}^2)\sum_jb_{n,j}^2\to0\), summation
gives
\[
 \log M_n(h)=\tfrac12\Var(S_n)h^2+o(1)
 \longrightarrow \tfrac12 v_\theta h^2,\qquad
 v_\theta=M_\theta''(0)>0.
\]
Thus \(M_\theta(h)=e^{v_\theta h^2/2}\), contrary to
\eqref{eq:theta-growth}. Alternatively, independence gives
\[
 \kappa_4(S_n)=\kappa_4(X)\sum_jb_{n,j}^4\longrightarrow0,
\]
contradicting the inherited strict fourth-cumulant sign in
\eqref{eq:moments}.
\end{proof}

\begin{corollary}[Weak convergence is also excluded]\label{cor:weak-pure}
The laws represented by \eqref{eq:pure-products} cannot converge weakly
to the normalized theta source.
\end{corollary}
\begin{proof}
Put \(v=\E X^2\). Integration by parts gives
\(\E X^4=1/4\), \(\E X^6=3v/4\). Strict Cauchy--Schwarz gives
\(1/16<3v^2/4\), so \(\kappa_4(X)=1/4-3v^2<0\).
Independence therefore gives
\(\E S_n^4\le3\Var(S_n)^2\). For positive variances,
Paley--Zygmund applied to \(S_n^2\)
shows
\(\mathbb P\{S_n^2>\Var(S_n)/2\}\ge1/12\).
Tightness forces bounded variances.

The function \(\log m(h)/h^2\) extends continuously at zero, is bounded
on finite intervals and tends to zero at infinity by
\eqref{eq:quartic-growth}. Hence \(m(h)\le e^{Ch^2}\) for all real \(h\).
The variance bound supplies a common sub-Gaussian bound for all \(M_n\).
It gives uniform integrability of every fixed exponential observable.
Weak convergence would thus give \(M_n(h)\to M_\theta(h)\) for real
\(h\). On any complex disc,
\(|M_n(z)|\le M_n(|z|)\), so normality and the identity theorem upgrade
this to complex-compact convergence, contradicting Theorem \ref{thm:pure}.
\end{proof}

The potential cone in Section \ref{sec:finite} is not the reason for this
exclusion. Indeed, for independent \(X,Y\) of the pure-quartic law,
the source of \(X+Y/4\) has a positive fourth and a negative sixth
potential coefficient. A finite exact calculation makes this explicit.
For general \(\delta>0\), put \(r=\delta^4\), \(A=1+r\), and
\(v=\Gamma(3/4)/\Gamma(1/4)\). Expanding
\(\int e^{-(s-\delta y)^4-y^4}\dd y\) at \(s=0\) gives
\begin{align*}
 [s^4](-\log\rho_\delta)
  &=\frac{1-(11/2)r+r^2+18rv^2}{A^3},\\
 [s^6](-\log\rho_\delta)
  &=\frac{8\delta^2v}{A^{9/2}}
       \bigl(-(1-r)^2+9rv^2\bigr).
\end{align*}
Differentiation is justified by quartic decay. The finite expansion uses
\(\E y^{4k}=(1/4)_k/A^k\) and
\(\E y^{4k+2}=(v/\sqrt A)(3/4)_k/A^k\).
For \(\delta=1/4\), strict Cauchy bounds \(1/12<v^2<1/4\) and exact
rational arithmetic \cite{records} give
\[
 0.97<[s^4](-\log\rho_{1/4})<0.99,\qquad
 -0.25<[s^6](-\log\rho_{1/4})<-0.11.
\]
Thus leaving the local potential cone can occur even without
inter-block interaction; Theorem \ref{thm:pure} supplies the stronger
global approximation obstruction.

The result concerns fixed pure-quartic elementary shapes. It does not
exclude variable quadratic tilts, dependent blocks, or a sequence with
changing elementary distributions outside this class. Its proof uses
field growth and cumulants, not an assumed location of theta zeros.

\section{Finite obstructions for other seed classes}\label{sec:finite}

\subsection{Independent Bernoulli products}
For real \(b_j\), every zero \(r\) of a finite product
\(\prod_j\cos(b_jz)\) forces a zero at \(qr\), for every odd integer \(q\).
The same holds after multiplication by a zero-free entire function.
The property is closed under nonzero complex-compact limits: Rouch\'e's
theorem gives approximant zeros \(r_n\to r\), and compact convergence at
\(qr_n\) then forces a limiting zero at \(qr\).

The complete-field Acb certificate \cite{records} gives, with
\[
 a=28.26945,\quad b=28.26946,\quad B=[3a,3b],
\]
\begin{center}
\begin{tabular}{ll}
\toprule
Input & Outward rational interval for \(H_0\)\\
\midrule
\(a\) & \((2\cdot10^{-11},3\cdot10^{-11})\)\\
\(b\) & \((-9\cdot10^{-10},-8\cdot10^{-10})\)\\
Every real \(x\in B\) & \((-3.68\cdot10^{-13},-3.67\cdot10^{-13})\)\\
\bottomrule
\end{tabular}
\end{center}
Thus \(H_0\) has a zero \(r\in(a,b)\) and is nonzero at \(3r\), without
any assumption about simplicity or the first zero.

\begin{proposition}\label{prop:cosine}
No independent Bernoulli cosine products, even with arbitrary counts and
zero-free entire multipliers, converge complex-compactly to \(H_0\).
Moreover, on \(K=[a,b]\cup B\), every real continuous function
\(f(x)=g(x)\prod_j\cos(b_jx)\) with \(g\ne0\) on \(K\) satisfies
\begin{equation}\label{eq:cosine-floor}
 \norm{f-H_0}_{C(K)}\ge3.6\cdot10^{-13}.
\end{equation}
\end{proposition}
\begin{proof}
The qualitative assertion follows from odd dilation and the certificate.
For the quantitative assertion, an error strictly below the stated
constant would preserve the opposite signs at \(a,b\). The intermediate
value theorem supplies a zero \(r_f\in(a,b)\). A cosine factor must vanish
there, and hence at \(3r_f\), contradicting the certified margin on \(B\).
\end{proof}

\subsection{A source-potential inequality}
For a positive source \(\rho\), write \(W(x)=-\log\rho(\sqrt x)\), and
let \(\Delta_h\) denote a forward difference in \(x\).
If
\begin{equation}\label{eq:potential-cone}
 -\log\rho(u)=c+\beta u^2+\alpha u^4+\sum_{k\ge3}v_ku^{2k},
 \qquad v_k\ge0,
\end{equation}
with convergence at the tested nodes, then \(\Delta_h^3W(0)\ge0\).
Terms of degree at most two in \(x\) cancel, and higher monomials have
positive third differences. Normalization and the quadratic heat weight
do not change the test. The complete-source certificate gives
\begin{equation}\label{eq:delta3}
 \Delta_{1/100}^3[-\log\Phi(\sqrt x)]_{x=0}
 \in(-0.000434999,-0.000434997).
\end{equation}
Consequently sources in \eqref{eq:potential-cone} cannot converge to theta
even at these four source nodes. We do not infer a general weak-limit
exclusion for a growing-degree cone from this pointwise test alone.

\begin{proposition}[Single quartic laws]\label{prop:single}
Normalized laws proportional to
\(e^{-\alpha_nu^4-\beta_nu^2}\), \(\alpha_n>0,\beta_n\in\R\), cannot
converge weakly to the normalized theta source.
\end{proposition}
\begin{proof}
If \(r_n=\norm{(\alpha_n,\beta_n)}\to\infty\), pass to a subsequence of
normalized coefficient vectors with limit \((a,b)\) of norm one.
The polynomial \(p(u)=au^4+bu^2\) is nonconstant. Choose bounded intervals
\(A,B\) with \(\inf_Ap>\sup_Bp+2\eta\). Uniform convergence of the
normalized potentials on these intervals gives
\[
 \mu_n(A)\le(|A|/|B|)e^{-r_n\eta}\mu_n(B)\longrightarrow0,
\]
contradicting positive theta mass on \(A\).
Thus the parameters are bounded. Extract their limits. The normalizer
is determined by the mass of any compact interval of positive limiting
mass, so local density convergence identifies the limit as a quartic
or Gaussian law. Its third difference is zero, contrary to
\eqref{eq:delta3}. Integrability rules out a nonintegrable boundary
potential. Normalized Fourier convergence would imply weak convergence
by L\'evy's continuity theorem and is excluded as well.
\end{proof}

\subsection{Fixed Bessel base with positive even coefficients}
Let
\begin{equation}\label{eq:bessel}
 B_0(u)=e^{-2\pi\cosh4u},\qquad
 \int_0^\infty B_0(u)\cos(zu)\dd u=\tfrac14K_{iz/4}(2\pi).
\end{equation}
The Bessel seed has only real zeros; see \cite{gasper,dlmf}.
Even universal multipliers \(U\), with \(U(iz)\in\LP\) and
\(U\ge0\) on \(\R\), preserve this property \cite[Proposition 1]{cardon}.
Their canonical products have the form
\[
 U(u)=Cu^{2m}e^{au^2}\prod_j(1+u^2/r_j^2),
 \qquad C>0,\quad a\ge0.
\]
In particular all even Taylor coefficients are nonnegative.
The following larger-class exclusion does not require universality.

\begin{proposition}\label{prop:bessel}
Suppose \(U_n(u)=\sum_{k\ge0}c_{n,k}u^{2k}\) are nonzero even entire
functions with \(c_{n,k}\ge0\), and \(B_0U_n\) are integrable nonnegative
sources. Their even
source measures cannot converge weakly to the theta source, nor can their
cosine transforms converge to \(H_0\) on every complex compact.
\end{proposition}
\begin{proof}
Every \(U_n(\sqrt x)\) has nonnegative fourth forward differences.
The complete-source certificate instead gives
\begin{equation}\label{eq:delta4}
 \Delta_{1/1000}^4
 [\Phi(\sqrt x)e^{2\pi\cosh(4\sqrt x)}]_{x=0}
 \in(-5.232\cdot10^{-7},-5.230\cdot10^{-7}).
\end{equation}
For the weak-limit assertion, bounded masses imply, for every \(R\ge0\),
\[
 U_n(R)\int_R^{R+1}B_0(u)\dd u
 \le\int_0^\infty B_0(u)U_n(u)\dd u\le C.
\]
Since \(|U_n(z)|\le U_n(|z|)\), the multipliers form an entire normal
family. A compactly convergent subsequence gives local uniform density
convergence to \(B_0U\), which weak convergence identifies with \(\Phi\).
This contradicts \eqref{eq:delta4}. Fourier convergence implies bounded
masses and, after normalization, weak convergence by L\'evy's theorem.
\end{proof}

De Bruijn already discusses the universal-factor approximation obstruction
for this fixed Bessel base \cite[p. 221]{debruijn}. Equation
\eqref{eq:delta4} is a finite diagnostic of the stated coefficient class,
not a priority claim for that classical obstruction. Varying or
degenerating Bessel bases and multipliers outside the coefficient cone
remain outside Proposition \ref{prop:bessel}.

\section{Admissible coupled seeds and their precise limitations}
\label{sec:construction}

\subsection{Classical stability and exact heat evolution}
For \eqref{eq:network}, the continuous-spin Lee--Yang theorem gives
nonvanishing of the Laplace transform when all site fields have positive
real part \cite[Theorems 1 and 3]{sg}; see also \cite[Theorem 1]{fr}.
Every site is an even quartic law, the couplings are nonnegative, and
positive readout weights put all site fields in the same half-plane.
Spin-flip symmetry gives the negative half-plane as well.
The transform \(M(iz)\) therefore has only real zeros.
The same fixed law, weighted by \(e^{tS^2}\), has
\begin{equation}\label{eq:continuous-heat}
 \beta_j(t)=\beta_j-tb_j^2,\qquad
 K_{ij}(t)=K_{ij}+2tb_ib_j,\qquad
 \partial_t Z=\partial_h^2 Z,
\end{equation}
where \(Z(t,h)=\E_\nu e^{tS^2+hS}\) retains its initial normalizer.
All finite complex derivatives converge on compact sets by quartic
domination. The real-zero conclusion is asserted for real \(t\ge0\).

\subsection{An explicit seed outside the potential cone}
Take
\begin{equation}\label{eq:two-seed}
 \dd\nu(x,y)=C e^{-x^4-y^4+(7/2)(x^2+y^2)+xy}\dd x\dd y,
 \qquad S=x+y.
\end{equation}
It is initially interacting. Its heat parameters are
\(\beta(t)=-7/2-t\) and \(K(t)=1+2t\).
For \(s=x+y\), \(v=(x-y)/2\), its marginal is proportional to
\[
 e^{-s^4/8+2s^2}
 \int_\R e^{-2v^4+6v^2-3s^2v^2}\dd v.
\]
Let \(W=V^2\) under the normalized law
\(V\propto e^{-2(V^2-3/2)^2}\). Then
\begin{equation}\label{eq:sixth}
 [s^6](-\log\rho_S(s))=\tfrac92\kappa_3(W).
\end{equation}
The outward integral certificate \cite{records} encloses this coefficient
by
\[
 -0.02276664341391261483790292683089089842
 \quad\text{with radius }7.60\cdot10^{-39}.
\]
Thus this marginal is outside \eqref{eq:potential-cone}. It invalidates
an extension of that local cone exclusion to all ferromagnetic marginals.
It is nevertheless not a theta source: Theorem \ref{thm:persistent}
excludes this fixed seed, and also appending sites while retaining one
of its fixed diffuse readout sites.

\subsection{Elementary Ising realization, with a paid convergence domain}
The construction is a specialization of the classical quartic Ising
approximation \cite{sg}. For two even-size blocks of \(N\) unit spins,
let \(M_1,M_2\) be their magnetizations and set
\begin{equation}\label{eq:ising}
 c_N=(12N^3)^{-1/4},\qquad
 J_{\rm within}=1/N+7c_N^2,\qquad
 J_{\rm between}=c_N^2,\qquad
 S_N=c_N(M_1+M_2).
\end{equation}
Each \(J\) is the coefficient of an unordered spin pair.
Every coupling is positive at time zero.

We spell out the convergence payment for this exact normalization.
For even \(N\) and \(M=2r\), the centered binomial ratio is
\[
 R_N(M)=\frac{\binom N{(N+M)/2}}{\binom N{N/2}}
 =\prod_{j=1}^{|r|}
 \frac{1-(2j-1)/(N+1)}{1+(2j-1)/(N+1)}.
\]
The inequality \(\log[(1-u)/(1+u)]\le-2u-2u^3/3\) gives
\begin{equation}\label{eq:binomial}
 \log R_N(M)\le-\frac{M^2}{2(N+1)}
  -\frac{M^2(M^2-2)}{12(N+1)^3}.
\end{equation}
Writing \(x=c_NM\), the within-block weight therefore has the uniform
central-ratio envelope
\begin{equation}\label{eq:envelope}
 R_N(M)e^{J_{\rm within}M^2/2}
 \le e^{-x^4/8+(7/2+\sqrt{12})x^2}.
\end{equation}
On fixed \(x\)-compacts, expansion of the same exact ratio gives
\[
 \log R_N(M)=-\frac{M^2}{2N}-\frac{M^4}{12N^3}+o(1);
\]
the error is uniform there, since \(M=O(N^{3/4})\).
Thus the central ratio tends to \(e^{-x^4+(7/2)x^2}\).
The cross-block factor is \(e^{xy}\).

Use associated piecewise-constant histograms of mesh \(2c_N\).
Inequality \eqref{eq:envelope}, with slightly weaker constants inside
each shrinking cell, gives an integrable quartic envelope.
The cross term is absorbed by \(|xy|\le(x^2+y^2)/2\).
The local limits and dominated convergence give convergence of
normalizers and densities to \eqref{eq:two-seed}. The actual lattice
readout differs from the histogram readout by at most the shrinking
mesh. On each compact complex \((t,h)\) set the same quartic envelope
absorbs \(e^{t(x+y)^2+h(x+y)}\) and each fixed derivative power.
The normalized heat transforms and every fixed derivative consequently
converge complex-compactly to the specified continuous seed.

At finite \(N\), heat adds \(2tc_N^2\) to every pair and retains the
diagonal scalar \(e^{2tNc_N^2}\), with the initial normalizer fixed.
This gives \eqref{eq:continuous-heat} exactly before taking the limit.
This argument is for a fixed two-block architecture and its explicit
continuous limit; it gives no theta approximation or estimate uniform
in an increasing number of blocks.

\section{Certificates, controls and research disposition}\label{sec:records}

The replay records contain small outward enclosures rather than large
generated data. Their source hashes identify the implementations and do
not replace the arithmetic or the analytic implications above.
The dependencies used were Python 3.10.0, python-flint 0.9.0 and
FLINT 3.6.0; the arithmetic contracts are documented in \cite{flint}.

\begin{center}\small
\begin{tabularx}{\textwidth}{@{}p{.29\textwidth}X@{}}
\toprule
Record / evaluator & Fixed retained run and scope\\
\midrule
Harmonic sign certificate &
160 bits, three complete xi evaluations, a 30-second cap.
One ball covers the entire tripled interval; no source cutoff.\\
Source alternative screens &
160 bits, nine source evaluations with eight theta terms and a proved
complete tail; 30-second cap and required difference width below \(10^{-7}\).\\
Two-quartic coefficient &
160 bits, four Acb integrals on \([0,4]\); goals \(2^{-130}\), at most
100,000 evaluations and depth 100 per integral.\\
Unequal quartic convolution &
Exact rational enumeration through degree six at \(\delta=1/4\);
three coefficients, no quadrature or parameter grid.\\
\bottomrule
\end{tabularx}
\end{center}

For the source screens, positivity and geometric summation give
\[
 0<T_N(u)\le
 \frac{2\pi^2e^{9u}(N+1)^4e^{-\pi(N+1)^2}}
 {1-16e^{-\pi(2N+3)}}\qquad(u\ge0).
\]
For the four integrals defining \(W\), the omitted tails satisfy, for
\(u\ge4\), \(0\le k\le3\),
\[
 u^{2k}e^{-2(u^2-3/2)^2}\le e^{-u^4/2}\le e^{-32u},
 \qquad \int_4^\infty e^{-32u}\dd u=e^{-128}/32.
\]
Both intervals and quadrature contracts are therefore paid independently
of agreement between floating approximations. Scouts preceded the
retained certificate contracts and are identified as such in the records.
The new field-growth results in Sections \ref{sec:persistent} and
\ref{sec:indquartic} require no additional numerical campaign.

The repository's collision controls \cite{controls} screen any proposed
general theta implication. A claim of stable complex-compact approximation
to an Euler-factor control with \(0<c<1\) fails at time zero, where its
explicit extra factor already has nonreal zeros. For the weighted-square,
reflected, and smooth/positive-primitive positive-threshold controls,
the corresponding convergence fails on a positive pre-collision slice
with a nonreal zero. Hurwitz gives the contradiction on that slice.
A double real zero at the collision time itself is compatible with
Lee--Yang. We have not assigned an uncomputed coefficient or parameter
representation to any control.

The results narrow the constructive search in concrete ways. Independent
Bernoulli products and independent scaled pure-quartic blocks cannot give
the desired limit, even with changing counts. Single quartic laws and
the specified fixed Bessel multiplier class are also excluded. Interactions
do change zero geometry, and coupled quartic marginals can leave the
simple source-potential cone. Nevertheless, any successful quartic network
must satisfy the degeneration condition \eqref{eq:rectangle}; appending
sites while retaining a fixed diffuse quartic seed cannot suffice.

What remains is a construction in the infinitesimal or degenerating
interacting regime, with an actual theta approximation error tending to
zero and uniform heat/differentiation bounds. Neither the classical
stability theorem, the explicit seed, nor the exclusion results supply
that construction. The investigation remains open in that stated regime,
pending formal review of the present round.

\section*{Acknowledgement of LLM assistance}
This working manuscript was drafted for Edward Baker with substantial
assistance from OpenAI \texttt{gpt-6.1-sol} (Codex), reasoning effort
\texttt{ultra}. Same-model agents assisted with literature checks,
derivations, source arithmetic and authoring checks. These checks do not
constitute independent cross-model or specialist human review. The
manuscript is unreviewed and does not imply the intended author's
endorsement of its mathematical claims.

\begin{thebibliography}{99}
\bibitem{polymath}
D. H. J. Polymath,
\emph{Effective approximation of heat flow evolution of the Riemann
\(\xi\) function, and a new upper bound for the de Bruijn--Newman constant},
Research in the Mathematical Sciences 6 (2019), article 31.
\href{https://arxiv.org/abs/1904.12438}{arXiv:1904.12438}.
\bibitem{sg}
B. Simon and R. B. Griffiths,
\emph{The \((\phi^4)_2\) field theory as a classical Ising model},
Communications in Mathematical Physics 33 (1973), 145--164.
\url{https://math.caltech.edu/SimonPapers/47.pdf}.
\bibitem{fr}
J. Fr\"ohlich and P.-F. Rodriguez,
\emph{Some applications of the Lee--Yang theorem},
Journal of Mathematical Physics 53 (2012), 095218.
\href{https://arxiv.org/abs/1205.6643}{arXiv:1205.6643}.
\bibitem{debruijn}
N. G. de Bruijn,
\emph{The roots of trigonometric integrals},
Duke Mathematical Journal 17 (1950), 197--226.
\url{https://repository.tudelft.nl/file/File_524a4d3c-0858-4c55-b425-8c8ce1fb89da}.
\bibitem{cardon}
D. A. Cardon,
\emph{Fourier transforms having only real zeros}, Proposition 1.
\url{https://mathdept.byu.edu/cardon/papers/FourierRealZeros.pdf}.
\bibitem{gasper}
G. Gasper,
\emph{Using integrals of squares of certain real-valued special functions
to prove that the P\'olya \(\Xi^*(z)\) function, the functions \(K_{iz}(a)\),
\(a>0\), and some other entire functions have only real zeros}, 2008.
\href{https://arxiv.org/abs/0801.2996}{arXiv:0801.2996}.
\bibitem{dlmf}
NIST Digital Library of Mathematical Functions,
Sections 5.11, 10.25, 10.40 and 25.4, accessed 17 October 2026.
\url{https://dlmf.nist.gov/}.
\bibitem{cgn}
F. Camia, C. Garban and C. M. Newman,
\emph{Planar Ising magnetization field II. Properties of the critical and
near-critical scaling limits},
\href{https://arxiv.org/abs/1307.3926}{arXiv:1307.3926}.
\bibitem{flint}
The python-flint authors,
\emph{Arb and Acb numerical evaluation and analytic integration documentation}.
\url{https://python-flint.readthedocs.io/en/latest/}.
\bibitem{scout}
Shifted zeta positivity research repository,
\emph{Theta moments exclude two spin families and screen a third},
10 October 2026, internal research note and outward Decimal record.
Repository path:
\path{papers/quasi-rh/quasi-rh-exponent-descent/newman_collisions/15_lee_yang/}.
\bibitem{records}
Shifted zeta positivity research repository,
\emph{Constructive stable models, round-one Notes 01--03 and certificates},
17 October 2026, unreviewed internal research records.
Repository paths:
\path{papers/investigations-3/05_constructive_stable_models/notes/round-1/}
and
\path{papers/investigations-3/05_constructive_stable_models/numerics/round-1/}.
\bibitem{controls}
Shifted zeta positivity research repository,
\emph{Collision controls and the screening they impose},
\path{result-summaries/12_collision_controls_and_screening.md},
with canonical source notes cited there.
\end{thebibliography}
\end{document}
